% ===== PRÉAMBULE CANONIQUE — à coller VERBATIM en tête de CHAQUE .tex =====
\documentclass[11pt,a4paper]{article}
\usepackage[utf8]{inputenc}\usepackage[T1]{fontenc}\usepackage{lmodern}
\usepackage[french]{babel}
\usepackage{amsmath,amssymb,mathtools}\usepackage{icomma}
\usepackage{siunitx}\sisetup{locale=FR}  % \qty{}{}, \si{}, \ang{}
\usepackage{xcolor}\usepackage{colortbl}
\usepackage{tikz}\usepackage{pgfplots}\pgfplotsset{compat=1.18}
\usepackage[siunitx,europeanresistors]{circuitikz} % circuits (élec.)
\usepackage[most]{tcolorbox}\usepackage{enumitem}\usepackage{array,booktabs}
\usepackage[a4paper,margin=2.2cm]{geometry}
\usetikzlibrary{arrows.meta,calc,angles,quotes,positioning,patterns,%
  backgrounds,shapes.geometric,decorations.pathmorphing,decorations.markings,intersections}
\definecolor{primary}{RGB}{222,49,99}\definecolor{primarydark}{RGB}{150,22,55}
\definecolor{defcol}{RGB}{0,114,178}\definecolor{methcol}{RGB}{52,92,148}
\definecolor{excol}{RGB}{0,135,90}\definecolor{attncol}{RGB}{200,40,55}
\tcbset{boxbase/.style={enhanced,breakable,boxrule=0.9pt,arc=2.5pt,
  left=3mm,right=3mm,top=2.3mm,bottom=2.3mm,fonttitle=\bfseries,coltitle=white}}
% --- boîtes TUTORIEL (noms EXACTS, ne pas renommer) ---
\newtcolorbox{cle}[1][]{boxbase,colback=defcol!6,colframe=defcol,
  title={\textsc{À retenir}\ifx\relax#1\relax\relax\else\textmd{ : \itshape #1}\fi}}
\newtcolorbox{methode}[1][]{boxbase,colback=methcol!7,colframe=methcol,sharp corners,
  title={\textsc{Méthode}\ifx\relax#1\relax\relax\else\textmd{ --- \itshape #1}\fi}}
\newtcolorbox{exemplebox}[1][]{boxbase,colback=excol!5,colframe=excol,
  title={\textsc{Exemple}\ifx\relax#1\relax\relax\else\textmd{ : \itshape #1}\fi}}
\newtcolorbox{piege}[1][]{boxbase,colback=attncol!6,colframe=attncol,
  title={\textsc{Piège}\ifx\relax#1\relax\relax\else\textmd{ : \itshape #1}\fi}}
\newtcolorbox{remarque}[1][]{boxbase,colback=black!4,colframe=black!45,
  title={\textsc{Remarque}\ifx\relax#1\relax\relax\else\textmd{ : \itshape #1}\fi}}
% --- boîtes EXERCICE / CORRIGÉ (noms EXACTS, auto-comptées) ---
\newtcolorbox[auto counter]{exo}[1][]{boxbase,colback=primary!4,colframe=primary,
  title={\textsc{Exercice}~\thetcbcounter\ifx\relax#1\relax\relax\else\textmd{ --- \itshape #1}\fi}}
\newtcolorbox[auto counter]{corrige}[1][]{boxbase,colback=excol!4,colframe=excol,
  title={\textsc{Corrigé}~\thetcbcounter\ifx\relax#1\relax\relax\else\textmd{ --- \itshape #1}\fi}}
\newcommand{\vv}[1]{\vec{#1}}\newcommand{\ds}{\displaystyle}
\newcommand{\R}{\mathbb{R}}\newcommand{\N}{\mathbb{N}}\newcommand{\Z}{\mathbb{Z}}
% (un \usepackage{fancyhdr}+en-tête est possible ; il est ignoré côté web)
% =========================================================================
\begin{document}
\begin{exo}[objet suspendu : deux montages]
Un objet de masse $m=\qty{12}{kg}$ est suspendu par \textbf{deux fils} à deux dynamomètres. On
compare un montage \emph{symétrique} (cas 1) et un montage \emph{asymétrique} (cas 2). Le système
reste en équilibre, $g=\qty{10}{N/kg}$.
\begin{center}
\begin{tikzpicture}[>=Stealth,scale=0.9]
  \begin{scope}
    \draw[black!55,line width=1pt](-1.4,2)--(1.4,2);
    \foreach \x in {-1.2,-0.8,...,1.3}\draw[black!45](\x,2)--(\x-0.18,2.2);
    \coordinate(o)at(0,0.3);
    \draw[defcol,line width=1pt](-1,2)--(o);\draw[defcol,line width=1pt](1,2)--(o);
    \draw[fill=primary!25,draw=primary!70](-0.3,-0.2)rectangle(0.3,0.3);
    \draw[->,line width=1pt,orange](0,-0.2)--(0,-1) node[below]{$\vv G$};
    \node at (0,-1.7){\footnotesize Cas 1 (symétrique)};
  \end{scope}
  \begin{scope}[shift={(5,0)}]
    \draw[black!55,line width=1pt](-1.6,2)--(1.6,2);
    \foreach \x in {-1.4,-1,...,1.5}\draw[black!45](\x,2)--(\x-0.18,2.2);
    \coordinate(o)at(0,0.3);
    \draw[defcol,line width=1pt](-0.3,2)--(o);\draw[attncol,line width=1pt](1.4,2)--(o);
    \draw[fill=primary!25,draw=primary!70](-0.3,-0.2)rectangle(0.3,0.3);
    \draw[->,line width=1pt,orange](0,-0.2)--(0,-1) node[below]{$\vv G$};
    \node at (0,-1.7){\footnotesize Cas 2 (asymétrique)};
  \end{scope}
\end{tikzpicture}
\end{center}
Pour chaque affirmation, réponds par \textbf{VRAI} ou \textbf{FAUX} et justifie.
\begin{enumerate}[label=\alph*)]
  \item Dans le cas 1, chaque dynamomètre indique \qty{60}{N}.
  \item La somme des contributions verticales des deux tensions vaut \qty{240}{N}.
  \item Dans le cas 2, les deux dynamomètres indiquent encore la même valeur.
  \item En passant du cas 1 au cas 2, la charge totale supportée par l'ensemble des fils change.
\end{enumerate}
\end{exo}

\begin{exo}[plaque sur deux tiges]
Une plaque rigide est suspendue par deux tiges $D_1$ et $D_2$ placées à ses extrémités. On dépose
un objet de masse \qty{20}{kg} au \emph{milieu} : chaque tige indique alors \qty{300}{N}
($g=\qty{10}{N/kg}$).
\begin{enumerate}[label=\alph*)]
  \item Quelle est l'intensité totale de la force d'attraction terrestre s'exerçant sur tout
        l'ensemble (objet + plaque + tiges) ?
  \item Un étudiant affirme : « la masse de la plaque vaut \qty{300}{kg} ». Montre par un calcul que
        c'est \emph{impossible}.
  \item On déplace l'objet vers $D_1$ jusqu'à ce que $D_1$ affiche \qty{380}{N}. Que lit alors $D_2$ ?
\end{enumerate}
\end{exo}

\begin{exo}[suspension par deux câbles inclinés]
Un objet de masse $m$ est suspendu par \textbf{deux câbles identiques} accrochés au plafond, chacun
faisant un angle $\theta$ avec la verticale (montage symétrique, $g=\qty{10}{N/kg}$).
\begin{center}
\begin{tikzpicture}[>=Stealth,scale=1]
  \draw[black!55,line width=1pt](-2,2)--(2,2);
  \foreach \x in {-1.8,-1.4,...,1.9}\draw[black!45](\x,2)--(\x-0.2,2.22);
  \coordinate(o)at(0,0.2);
  \draw[defcol,line width=1.1pt](-1.2,2)--(o) node[midway,left]{$\vv T$};
  \draw[defcol,line width=1.1pt](1.2,2)--(o) node[midway,right]{$\vv T$};
  \draw[black!45,densely dashed](0,2)--(0,0.2);
  \draw[->,black!60] (0,1.3) arc (-90:-124:0.7); \node at (-0.42,1.35){\scriptsize $\theta$};
  \draw[fill=primary!25,draw=primary!70](o) circle (0.2);
  \draw[->,line width=1.2pt,orange](0,0)--(0,-1) node[below]{$\vv G$};
\end{tikzpicture}
\end{center}
\begin{enumerate}[label=\alph*)]
  \item Fais le bilan des forces, puis projette l'équilibre sur la verticale pour établir que
        $2T\cos\theta = G$.
  \item Déduis-en l'expression de la tension $T$ d'un câble en fonction de $m$, $g$ et $\theta$.
  \item Application : $m=\qty{20}{kg}$ et $\theta=\ang{60}$. Calcule $T$ et compare-la à $G/2$.
  \item Chaque câble rompt au-delà de \qty{500}{N}. Quelle masse maximale peut-on suspendre à
        $\theta=\ang{60}$ ?
\end{enumerate}
\end{exo}

\begin{exo}[équilibre à vitesse constante]
On rappelle qu'un corps dont la vitesse est \textbf{constante} (mouvement rectiligne uniforme) a une
accélération nulle : il est en équilibre de translation, $\sum\vv F=\vv 0$.
\begin{enumerate}[label=\alph*)]
  \item Un parachutiste de masse $m=\qty{80}{kg}$ descend à vitesse constante ($g=\qty{10}{N/kg}$).
        Il n'est soumis qu'à son poids et à la force de frottement de l'air (verticale, vers le haut).
        Détermine l'intensité de cette force de frottement.
  \item Un avion vole en ligne droite à vitesse constante. Il subit quatre forces : son poids, la
        portance (verticale, vers le haut), la force motrice des réacteurs (horizontale, vers l'avant)
        et la résistance de l'air (horizontale). Indique quelles forces se compensent deux à deux, et
        avec quelle intensité relative.
\end{enumerate}
\end{exo}

\begin{exo}[chasse aux affirmations fausses]
Indique si chaque affirmation est \textbf{correcte} ou \textbf{fausse}, en corrigeant les fausses
($g=\qty{10}{N/kg}$).
\begin{enumerate}[label=\alph*)]
  \item Un corps de masse \qty{200}{g} a un poids de \qty{200}{N}.
  \item Lorsqu'un objet change d'endroit sur la Terre, c'est son \emph{poids} qui peut varier
        légèrement, pas sa \emph{masse}.
  \item Une caisse posée au sol, tirée horizontalement par une force de \qty{30}{N}, reste immobile :
        la force de frottement exercée par le sol vaut alors \qty{30}{N}.
  \item Un solide tel que $\sum\vv F=\vv 0$ est nécessairement en équilibre.
\end{enumerate}
Enfin : \emph{combien} de ces quatre affirmations sont correctes ?
\end{exo}

\end{document}
